%% ACM MM 2026 -- Author Response (Rebuttal), Paper #7367
%% By-reviewer, point-by-point, in order. Reviewer header runs in with the first
%% point; repeated concerns point to the first full reply (see Rx-y).

\documentclass[10pt,twocolumn]{article}
\usepackage[top=0.45in,bottom=0.4in,left=0.45in,right=0.45in]{geometry}
\usepackage{newtxtext,newtxmath}
\usepackage{enumitem}
\usepackage{microtype}
\usepackage{amsmath}
\usepackage{xcolor}
\usepackage{eso-pic}
\usepackage[switch]{lineno}
\definecolor{rvR1}{RGB}{0,75,190}
\definecolor{rvR2}{RGB}{190,95,0}
\definecolor{rvR3}{RGB}{0,140,85}
\definecolor{rvR4}{RGB}{190,35,45}
\definecolor{rvR5}{RGB}{125,55,175}

\newcommand{\rvcolor}{rvR1}
\setlength{\parindent}{0pt}
\setlength{\parskip}{0.7pt}
\setlength{\columnsep}{0.15in}
\pagestyle{empty}
\newcommand{\rebuttalheader}{ACM MM 2026 Submission \#7367. Author Rebuttal.}
\AtBeginDocument{%
  \AddToShipoutPictureFG*{%
    \AtPageUpperLeft{%
      \raisebox{-0.24in}{%
        \makebox[\paperwidth]{\scriptsize\bfseries\rebuttalheader}%
      }%
    }%
  }%
}
\newcommand{\pt}[1]{\textbf{\textcolor{\rvcolor}{#1}}\ }
\newcommand{\rv}[1]{\vspace{1.5pt}\pt{#1}}

\begin{document}
\small
\linenumbers

We thank all five reviewers for their careful reading and are encouraged that
they find the work worthwhile: ``sensible'' (27Jx), ``important''/``promising''
(D8hx), and ``a fair and interesting contribution'' with ``compelling''
matched-setting results (wGgF). We take the presentation problems seriously and
have completed a full language and formatting revision, verifying that no model,
experiment, or reported result is affected by these edits.

\rv{Reviewer 27Jx (\#1): R1-1: Presentation.} We sincerely apologize for the
draft-stage errors, and we are very grateful for the reviewer's careful reading. We have carried out a thorough revision:
resolved the dangling \texttt{Figure ??} to Fig.~1 (Sec.~3.2), corrected the flagged
typos (\texttt{multi-granul}$\to$\texttt{multi-granular}, \texttt{Constrion}$\to$\texttt{Construction},
\texttt{Fame-level}$\to$\texttt{Frame-level}), re-rendered Figs.~1 and 2 as vector graphics,
and made a careful, line-by-line language pass over the entire manuscript. These
were purely presentation issues, not substantive ones. No result, number, or claim was affected.

\pt{R1-2: Distinction from Vgent.} We thank the reviewer.
The distinction is paradigmatic and does not hinge on any single layer. Vgent is
a \emph{fixed} retrieve-then-read graph-RAG over frame and segment local
relations. In contrast, our contribution is twofold: we build a multi-level
scene-graph index over long video, and we design a modular, plug-and-play
adaptive agent that uses the scene graph as its primary retrieval index,
verifying each proposal against raw frames before answering (R2-3). Together they drive the +18.2-pp
matched-setting gain on LVBench (54.8\%$\to$73.0\%). We address the event layer's
specific role, and its accuracy contribution, in R1-3.

\pt{R1-3: The role of the event graph.} We 
agree with the observation. At a fixed backbone the event layer's measurable
effect is efficiency: it cuts LLM and frame-inspection calls by 71.5\% and 74.7\%
at unchanged accuracy (64.6\%$\to$64.4\%). We read
this as a capability gain, not a void: reaching the same answers with far less
search is itself better reasoning. Still, efficiency need not carry the
contribution alone. The flat
accuracy reflects the conservative event-graph construction of the submitted
version, not an intrinsic ceiling. In an additional experiment that strengthens
the construction and queries the index \emph{caption-free}, the agent reaches
\textbf{68.5\%} e2e accuracy on LVBench, a 3.9-pp gain over the segment-only
version and a 13.7-pp gain over the caption-only baseline. A graph that answers
68.5\% of questions with \emph{no} captions is an answer-bearing reasoning
substrate, so the distinction does \emph{not} collapse. The same structured index lets the lighter \texttt{o4-mini} reach
73.0\%, matching DVD's far stronger \texttt{o3} (74.2\%) and beating DVD's own
\texttt{o4-mini} (70.2\%). Our continued experiments improve these results
further, and we will report fair, matched comparisons in the revision.

\renewcommand{\rvcolor}{rvR2}%
\rv{Reviewer D8hx (\#2): R2-1: Offline construction cost.} For all 103 LVBench
videos (115.5h), construction costs \$112.74 (12,360 \texttt{gpt-4.1} frame-graph
calls at one keyframe per 5s, $\le$120/video, 7.94h on four workers, and the symbolic
Stage~3 event build adds only 121s), i.e.\ \$0.98 per video-hour. We reuse DVD's open-sourced captions rather than
generating our own, so this is our only added offline cost. End-to-end VQA over all 1549 questions costs \$279.54
(\texttt{o4-mini}) and \$499.31 (\texttt{o3}), so this one-time construction is
only 40\% and 23\% of the recurring online budget it serves.

\pt{R2-2: Cross-benchmark control.} We agree not all entries are controlled. On
LVBench the DVD comparison \emph{is} fully controlled: the same \texttt{o3}
backbone, DVD's own open-sourced captions, and the same video database. On the
other three, DVD released no full reproduction setup, so we cannot match its
pipeline. We still use its captions and report these as best-effort positions.
Other commercial and open models have closed weights or fixed pipelines, so they are reference
positions only. Our claims
therefore rest on the fully controlled LVBench comparison and its matched-setting
ablations.

\pt{R2-3: Direct graph quality.} Direct evaluation is desirable, but intrinsic
scene-graph metrics need a gold-annotated benchmark that does not exist at hour
scale, and they would still miss what matters: extraction accuracy and the
graph's net effect on answers. Comparable systems, DVD included, are judged the
same way, not by intrinsic quality. We therefore measure the graph through the
task: the caption-free graph-only setting is cleanest, since with captions
removed the answer must come from the graph, yet accuracy holds at 68.5\% on
LVBench (Please see R1-3), and the matched ablations isolate its gain. We will still add a
small human-verified sample for entity and relation precision. The pipeline
already shows 76.8\% entity-mention absorption and a 5.0\% invalid-relation rate.
The propose-and-verify re-inspection of the raw frames further reduces how much
frame-level error propagates into the event graph.

\pt{R2-4: Implementation details.} To give the values asked for, entity alignment
scores name, description, and type similarity with weights 0.7, 0.2, and 0.1, and
merges above a threshold of 0.78. The entity blacklist removes artefacts and
over-generic terms, and segments merge into events at a 12s temporal gap with
relation-overlap thresholds. The composite retrieval score and all remaining thresholds are fully specified in
the appendix, which gives the alignment, gating, and event-construction
algorithms. We will release a fully reproducible repository with the revision.

\pt{R2-5: Presentation.} Thanks. Please see R1-1.

\pt{R2-6: Isolating localization from FTT.} Our reported answered accuracy
(ans\_acc) separates the two: the caption-only baseline scores 65.7\% and
SGDAgent 74.2\% (+8.5\,pp), independent of the answer rate, while the FTT drop
(16.5\%$\to$1.6\%) is itself a consequence of structured retrieval converging
where open caption search loops for lack of locatable evidence.

\renewcommand{\rvcolor}{rvR3}%
\rv{Reviewer k56s (\#3): R3-1: Baseline data citation.} Thank you for catching
this. The correct DVD values are Video-MME 67.3 (we had mis-copied DVD's
LongVideoBench(Long) entry, 68.6) and EgoSchema 76.6, both matching DVD's
published numbers. The correction in fact
widens our Video-MME margin to 69.0 vs.\ 67.3. We will present a carefully re-verified version, with a more
general Fig.~1 example, in the revision.

\pt{R3-2: Related works.} Thank you. We will cite and discuss all three ([1] Mind
Palace, [2] the Gaussian-token work, [3] GraphVideoAgent) in the revision.
They use the graph as an in-context representation or online entity memory,
whereas we use it as the agent's primary retrieval index, gated to raw-frame
verification before answering (Please see R1-2), a distinct design.

\pt{R3-3: Offline construction cost.} Thanks. Please see R2-1.

\pt{R3-4: Error propagation \& robustness.} Two mechanisms limit it. First,
the propose-and-verify step gates every graph hit to raw-frame re-inspection before answering
(R2-3), so a hallucinated or mis-aligned node is caught at verification rather
than read out as an answer, reducing upward propagation. Second, the caption-free
graph-only system still reaches 68.5\% (Please see R1-3), evidence that the pipeline
tolerates imperfect frame-level extraction. To quantify this we will add a
robustness analysis that perturbs frame-graph quality and reports the accuracy
sensitivity, and we will discuss [4].

\renewcommand{\rvcolor}{rvR4}%
\rv{Reviewer SSkz (\#4): R4-1: Over-strong coverage claim.} We agree, and will revise
the wording. The method improves \emph{localization within the sampled frames}, and
we will no longer claim to address incomplete sampling coverage, since the graph
itself is built only from sparsely sampled keyframes.

\pt{R4-2: The role of the event graph. } In the submitted table the event graph is
not a fixed-backbone accuracy driver, and we present it as a higher-order layer
that reduces the agent's search burden. A strengthened, caption-free run, however,
shows it also carries accuracy, not only efficiency (Please see R1-3).

\pt{R4-3: Heuristic construction.} Please see R2-3 for robustness. We now fully specify
the alignment, merging, and event thresholds (Please see R2-4).

\pt{R4-4: Presentation.} Please see R1-1. We have already re-rendered Figures.

\pt{R4-5: Questions.} Please see R2-3 for scene-graph accuracy. Repositioning the event
graph as an efficiency layer is appropriate because it reduces search burden, but a strengthened
caption-free run shows it is also accuracy-bearing (Please see R1-3). Please see R2-1 for full pipeline cost
per video. For the sampling interval, we currently sample one frame every 5s and will report
2s/5s/10s sensitivity in the revision.

\renewcommand{\rvcolor}{rvR5}%
\rv{Reviewer wGgF (\#5): R5-1: HIG \& hierarchy novelty.} Thank you. We agree HIG
proposes a hierarchical video graph, and we will cite and contrast it. We will not
claim the hierarchy concept as novel. Our work differs in task: HIG targets video
scene-graph \emph{generation}, whereas we use a graph index to \emph{drive an
adaptive agent} for long-video QA (Please see R1-2). 

\pt{R5-2: Implementation details.} The alignment, retrieval, gating, controller-signal, and
event-construction rules are all specified operationally in the appendix as
pseudocode. Please see R2-4 for the headline values.

\pt{R5-3: Graph-depth ablation.} Please see R1-3. The segment graph drives
fixed-backbone accuracy, and the event graph reduces search burden and, under a
stronger construction, also contributes accuracy (caption-free 68.5\%).

\pt{R5-4: Efficiency evidence.} Beyond online call counts we now also report
dollar cost and construction time (Please see R2-1).

\pt{R5-5: Presentation.} Thanks. Please see R1-1.

\end{document}
